\subsection{Arithmetic cohomology over the τ-base}\label{endpoint-tau_base--arithmetic-cohomology-over-the-ux3c4-base}

\subsubsection{Scope and result}\label{endpoint-tau_base--scope-and-result}

This note changes the geometric base of the existing arithmetic theta model to a one-point absolute base with distinguished absorbing element τ. It does \textbf{not} start with a variety over a finite field and then relabel its coefficients.

The main calculation is the right adjoint to the actual global-section functor on this τ-based model. A projective resolution computes the arithmetic cohomology; a dual injective complex computes its right-adjoint object. These are then linked to the existing finite Mellin-jet quotients, including the full nilpotent jets and the original e/τ distinction.

This is an algebraic, finite-chart realization. No assertion is made that the chart model is an arithmetic scheme, that its topology already supplies all prime correspondences, that the right-adjoint object is an involutive Verdier duality on every sheaf, or that the calculation proves a weight estimate. Each operation used below has its stated category. In particular the algebraic dual is not an assertion about the continuous dual of the Schwartz quotient.

The supplied SplitZero results are the source for the scalar operations, support fibres and transports. The new constructions and proofs below are calculations in this note; they have not been Lean-formalized.

\subsubsection{1. An actual base whose structural zero is τ}\label{endpoint-tau_base--an-actual-base-whose-structural-zero-is-ux3c4}

Let

\[
\mathbf F_{1,\tau}=\{\tau,1\}
\]

be the pointed multiplicative monoid with identity 1 and absorbing element τ. As a blueprint, its only additional zero relation is that the empty sum represents τ. No additive equation \(1+1=1\) is imposed. Its spectrum \(\mathfrak b_\tau\) has the one prime \(\{\tau\}\).

For every nonzero commutative ring R, retain

\[
G(R)=\{\tau\}\sqcup\{r^\bullet:r\in R\},\qquad e_R=0_R^\bullet.
\tag{1}
\]

Let \(B_R\) be its associated blueprint: the original multiplicative pointed monoid, with all its original finite additive relations. In particular retain \([\tau]=0\), in agreement with the presentation correction in the supplied formalization note. There is a blueprint morphism

\[
\jmath_R:\mathbf F_{1,\tau}\longrightarrow B_R,
\qquad \tau\mapsto\tau,\quad1\mapsto1_R^\bullet.
\tag{2}
\]

On spectra it gives a genuine structural morphism

\[
a_R:\operatorname{Spec}B_R\longrightarrow\mathfrak b_\tau.
\tag{3}
\]

Every prime has inverse image \(\{\tau\}\) under (2). Thus every arithmetic point lies over the single base point, by a specific inverse-image operation on prime ideals.

The arithmetic quotient stays attached:

\[
p_R:G(R)\longrightarrow R,
\quad p_R(\tau)=0_R,\quad p_R(r^\bullet)=r.
\tag{4}
\]

For a domain,

\[
P_\tau=\{\tau\},\qquad P_0=\{\tau,e_R\},\qquad
\operatorname{Spec}R\xrightarrow{i_R}\operatorname{Spec}_{\rm sr}G(R),
\quad\mathfrak p\mapsto\mathfrak p\sqcup\{\tau\}.
\tag{5}
\]

Here \(i_R(\operatorname{Spec}R)=V(e_R)=\overline{\{P_0\}}\). These maps carry the source's infinite quotient when \(R=\mathbb Z\). They do not assign an additional finite norm to either bottom ideal.

There is a second, different-direction morphism in the same diagram:

\[
\operatorname{Spec}\mathbb B\xrightarrow{\operatorname{Spec}\chi_R}
\operatorname{Spec}_{\rm sr}G(R)\xrightarrow{a_R}\mathfrak b_\tau,
\quad
\chi_R(\tau)=0,\quad\chi_R(r^\bullet)=1.
\tag{6}
\]

The first arrow lands at \(P_\tau\). The original localization calculation is

\[
G(R)[e_R^{-1}]\cong\mathbb B.
\tag{7}
\]

Indeed invertibility and \(e_R^2=e_R\) force \(e_R=1\), and \(e_Rr^\bullet=e_R\) then forces every supported scalar to equal 1. External absence remains distinct. We use (3), not localization (7), for the change of base. The arithmetic structure and its quotient remain upstream of (6).

For \(R=\mathbb Z\), the source's infinite target is retained in the square

\[
\begin{array}{ccc}
G(\mathbb Z)&\xrightarrow{G(j)}&G(\mathbb C)\\
p_{\mathbb Z}\downarrow&&\downarrow p_{\mathbb C}\\
\mathbb Z&\hookrightarrow&\mathbb C.
\end{array}
\tag{8}
\]

Here \(j:\mathbb Z\hookrightarrow\mathbb C\) is the specified unital embedding, and \(G(j)\) is its split lift.

\subsubsection{2. The arithmetic theta chart space over that base}\label{endpoint-tau_base--the-arithmetic-theta-chart-space-over-that-base}

Fix \(E=\mathbb C\). Let P be the four-point poset

\[
+<\eta<\sigma,\qquad -<\eta<\sigma,
\tag{9}
\]

with + and − incomparable. Give P the upper-set topology. Its proper nonempty opens are

\[
\{\sigma\},\quad\{\eta,\sigma\},\quad
\{+,\eta,\sigma\},\quad\{-,\eta,\sigma\}.
\]

The structural sheaf \(\mathcal S\) has value \(G(E)\) on every open containing an arithmetic chart point and value \(\mathbb B\) on \(\{\sigma\}\). Restriction to \(\sigma\) is \(\chi_E\); other nontrivial restrictions are identities. Compatible supported local amplitudes glue, and compatibility at σ distinguishes an absent family from a supported one. This proves the sheaf assertion directly.

Together with (2), this defines a space with a sheaf of blueprints over \(\mathfrak b_\tau\). There is also the explicit structural morphism to \(\operatorname{Spec}_{\rm sr}G(E)\), sending \(+, -,\eta\) to \(P_0\) and σ to \(P_\tau\), followed by the morphism induced by \(G(j)\) to \(\operatorname{Spec}_{\rm sr}G(\mathbb Z)\). On arithmetic stalks its map is \(G(\mathbb Q)\to G(E)\); at σ it is \(\mathbb B\to\mathbb B\). Thus the infinite quotient is still in a commuting structural diagram. This particular chart model maps to the arithmetic generic point; its finite-prime geometry is not being claimed to cover the closed prime fibres. Arithmetic prime information in the present realization enters through the theta map and its Euler/explicit formula, not through extra chart points invented here.

The coefficient spaces are exactly

\[
V=\{\phi\in\mathcal S(\mathbb R):\phi\text{ even},\ \phi(0)=0,\ \widehat\phi(0)=0\},
\quad
\widehat\phi(\xi)=\int_{\mathbb R}\phi(x)e^{-2\pi i x\xi}\,dx,
\tag{10}
\]

and

\[
\mathscr B=\{F\in C^\infty(\mathbb R_{>0}):
\sup_{x>0}x^b|(-x\partial_x)^kF(x)|<\infty
\ \text{for every }b\in\mathbb Z,k\ge0\}.
\tag{11}
\]

Define

\[
\Theta\phi(x)=\sum_{n\in\mathbb Z\setminus\{0\}}\phi(nx),
\qquad JF(x)=x^{-1}F(1/x).
\tag{12}
\]

The retained Poisson identity gives

\[
J\Theta\psi=\Theta\widehat\psi.
\tag{13}
\]

The arithmetic coefficient sheaf is the covariant diagram

\[
\mathcal T_+=V,\quad\mathcal T_-=V,\quad
\mathcal T_\eta=\mathscr B,\quad\mathcal T_\sigma=0,
\quad r_+=\Theta,\quad r_-=J\Theta=\Theta\mathcal F.
\tag{14}
\]

This use of a zero vector space at σ will not identify the supported zero with τ. It is reconstructed as the support skeleton in the next paragraph.

\paragraph{All original two-leg support faces}\label{endpoint-tau_base--all-original-two-leg-support-faces}

Set \(L=\mathcal P(\{+,-\})\). For nonempty \(A\subseteq\{+,-\}\), define a sheaf \(\mathcal T_A\) by

\[
(\mathcal T_A)_+=\begin{cases}V&+\in A\\0&+\notin A,\end{cases}
\quad
(\mathcal T_A)_-=\begin{cases}V&-\in A\\0&-\notin A,\end{cases}
\quad
(\mathcal T_A)_\eta=\mathscr B,\quad
(\mathcal T_A)_\sigma=0.
\tag{15}
\]

The nonzero restriction maps are those in (14); put \(\mathcal T_\varnothing=0\). Transports for \(A\subseteq A'\) are coordinate inclusions at + and − and the identity at η. These form a strict diagram.

Reconstruct the supplied SplitZero semimodule sheaf

\[
\widetilde{\mathcal T}
=\{\tau\}\sqcup\bigsqcup_{A\ne\varnothing}\{A\}\times\mathcal T_A.
\tag{16}
\]

At σ its stalk is the semilattice L, with Boolean scalar action. At an arithmetic chart the scalar action is by \(G(E)\). Restrictions to σ are semilinear along \(\chi_E\). In particular \(e_E(A,v)=(A,0)\), which restricts to A at σ, while τ restricts to the empty label. The original distinction is present at the new base boundary.

The additive module in degree zero is exactly the original two-leg carrier \((V^\tau)^2\): a missing leg is absent, and a supported zero in that leg is present. Formula (15) records a missing leg as a zero coefficient space \textbf{with its separate mask A retained}. Thus no new support is inserted into the original missing leg.

\subsubsection{3. Compute global cohomology over the τ-base}\label{endpoint-tau_base--compute-global-cohomology-over-the-ux3c4-base}

We now use the explicitly declared linear-fibre category

\[
\mathcal A=\operatorname{Fun}(P,\operatorname{Vect}_E).
\]

This is a realization category of the marked blueprint space, not an asserted category of étale sheaves on a new arithmetic scheme. Its passage to the original supported sheaves is (15)--(16), with kernels and quotients reconstructed fibrewise.

For \(x\in P\), let \(P_x\) be the projective diagram

\[
P_x(y)=\begin{cases}E&x\le y\\0&\text{otherwise}.\end{cases}
\quad
\operatorname{Hom}_{\mathcal A}(P_x,F)=F_x.
\tag{17}
\]

The constant diagram \(\mathbf E\) has the exact projective resolution

\[
0\longrightarrow P_\eta
\xrightarrow{v\mapsto(v,-v)}P_+\oplus P_-
\xrightarrow{(u,v)\mapsto u+v}\mathbf E\longrightarrow0.
\tag{18}
\]

At + and − the last map is identity. At η and σ it is addition with kernel the displayed anti-diagonal. This checks every stalk and every arrow.

Consequently the base pushforward is

\[
\mathsf A_\tau(F):=R\Gamma(P,F)
\simeq[F_+\oplus F_-\xrightarrow{r_+-r_-}F_\eta],
\quad\text{degrees }0,1.
\tag{19}
\]

This is global cohomology over the base point, rather than restriction to the added generic point. The typed comparison to the latter is

\[
\operatorname{res}_\sigma:\mathsf A_\tau(F)\longrightarrow F_\sigma[0].
\tag{20}
\]

Two chain representatives of (20) in degree zero are \((x,y)\mapsto r_{+\sigma}x\) and \((x,y)\mapsto r_{-\sigma}y\). Their difference is

\[
r_{\eta\sigma}(r_+x-r_-y),
\]

so they are chain homotopic, with homotopy \(r_{\eta\sigma}:F_\eta\to F_\sigma\). Thus (20) is canonical in the derived category.

On the actual joint arithmetic sheaf,

\[
\mathsf A_\tau(\mathcal T)
=[V\oplus V\xrightarrow{\Theta(\phi-\widehat\psi)}\mathscr B],
\tag{21}
\]

\[
H^0\mathsf A_\tau(\mathcal T)\cong V,
\quad \psi\mapsto(\widehat\psi,\psi),
\qquad
H^1\mathsf A_\tau(\mathcal T)=Q:=\mathscr B/\Theta V.
\tag{22}
\]

Injectivity of Θ follows from its Mellin transform on \(\Re s>1\), the nonvanishing Euler product there, and Mellin uniqueness. This is the existing source proof, not an assumption about zeros in the critical strip.

The actual supported differential is

\[
\widetilde d(\phi,\psi)
=(A,\Theta p(\phi)-J\Theta p(\psi)),
\tag{23}
\]

where A is the set of present legs; the absent pair maps to τ. Its cohomology uses

\[
Z^i_{\rm sp}=\operatorname{Eq}(d,\epsilon d),
\qquad
H^i_{\rm sp}=\operatorname{coeq}(b,\epsilon b),
\quad\epsilon(x)=ex.
\tag{24}
\]

For every nonempty A, \(H^1(\mathcal T_A)=Q\). For a single-leg A, \(H^0=0\); for the joint A, \(H^0=V\). Therefore

\[
\widetilde H^1_\tau
=\{\tau\}\sqcup\bigsqcup_{A\ne\varnothing}\{A\}\times Q.
\tag{25}
\]

Restriction (20) sends each of these arithmetic classes to the supported zero at its own label in the σ skeleton. Their global cohomology classes in (25) remain present. This is the exact answer to what remains when all geometric points lie over the τ-base.

\subsubsection{4. Compute the right adjoint on this base}\label{endpoint-tau_base--compute-the-right-adjoint-on-this-base}

Deligne defines a dualizing object through a right adjoint and then proves weight bounds. Here we compute the right adjoint to our stated functor (19), rather than importing a dualizing complex from a finite-field variety.

For \(x\in P\) and a vector space W, define

\[
I_x(W)(y)=\begin{cases}W&y\le x\\0&\text{otherwise}.\end{cases}
\quad
\operatorname{Hom}_{\mathcal A}(F,I_x(W))
=\operatorname{Hom}_E(F_x,W).
\tag{26}
\]

Because evaluation at x is exact and W is an injective E-module, \(I_x(W)\) is injective. The right-adjoint object is

\[
\boxed{
K_\tau(W)=
[I_\eta(W)\xrightarrow{(+\mathrm{res},-\mathrm{res})}
 I_+(W)\oplus I_-(W)],
\quad\text{degrees }-1,0.
}
\tag{27}
\]

For every sheaf F,

\[
\begin{aligned}
R\operatorname{Hom}_{\mathcal A}(F,K_\tau(W))
&=[\operatorname{Hom}(F_\eta,W)
 \longrightarrow\operatorname{Hom}(F_+,W)\oplus\operatorname{Hom}(F_-,W)]\\
&\cong R\operatorname{Hom}_E(\mathsf A_\tau(F),W).
\end{aligned}
\tag{28}
\]

The differential is \(\ell\mapsto(\ell r_+,-\ell r_-)\). This proves the adjunction, with its signs. For a bounded complex W, totalize the same formula with its original differential.

Let \(S_\eta W\) denote the diagram with W at η and zero at the other three points. Calculation of the four stalk complexes gives

\[
K_\tau(W)_+=[W\xrightarrow{1}W],\quad
K_\tau(W)_-=[W\xrightarrow{-1}W],\quad
K_\tau(W)_\eta=W[1],\quad K_\tau(W)_\sigma=0.
\tag{29}
\]

Thus

\[
\boxed{K_\tau(W)\simeq S_\eta W[1].}
\tag{30}
\]

The symbol \([1]\) means W is in cohomological degree −1. It is obtained from (18) and (27). The functor is a right adjoint to \(R\Gamma\) on this model; no assertion that it is the extraordinary pullback for an unconstructed proper arithmetic morphism is made. In particular, an involutive local Verdier duality or a finite-field purity theorem has not been obtained just from (28).

Apply (27) componentwise to support diagrams. Every zero coefficient stalk reconstructs as its labelled zero, including at σ. Dual transports are actual precomposition maps

\[
\rho_{AB}^{\vee}:\operatorname{Hom}(H_B,W)\longrightarrow
\operatorname{Hom}(H_A,W),\quad\ell\mapsto\ell\rho_{AB}.
\tag{31}
\]

They therefore have index category \(L^{\mathrm{op}}\); the label arrows have not been reversed without retaining the precomposition operation.

\subsubsection{5. Retain the actual characteristic-zero arithmetic action}\label{endpoint-tau_base--retain-the-actual-characteristic-zero-arithmetic-action}

For a positive real a, use

\[
U_aF(x)=F(x/a),\qquad
\mathcal M F(s)=\int_0^\infty F(x)x^s\,\frac{dx}{x}.
\tag{32}
\]

Then

\[
\mathcal M U_aF(s)=a^s\mathcal MF(s),\qquad
\widehat{U_a\phi}=aU_{1/a}\widehat\phi.
\tag{33}
\]

The action on the two source charts is consequently

\[
(\phi,\psi)\longmapsto(U_a\phi,\ aU_{1/a}\psi),
\quad
F\longmapsto U_aF\text{ at }\eta.
\tag{34}
\]

Using the same \(U_a\) on both source charts would fail the required intertwining. Direct calculation gives

\[
U_aJ=aJU_{1/a},\qquad d(U_a\phi,aU_{1/a}\psi)=U_a d(\phi,\psi).
\tag{35}
\]

It fixes every support mask. At σ it fixes the support skeleton. The new base remains the same point; the induced action on its pushforward remains the arithmetic action (33).

The source's moment map gives the two coefficient lines

\[
\mathcal L_0=E\text{ with action }1,\quad
\mathcal L_1=E\text{ with action }a,
\quad
m(\phi)=(\phi(0),\int_{\mathbb R}\phi(x)\,dx),
\tag{36}
\]

\[
mU_a=\operatorname{diag}(1,a)m.
\]

Both choices \(K_\tau(\mathcal L_0)\) and \(K_\tau(\mathcal L_1)\) exist. The latter produces, on dual functionals, the action

\[
(U_a\ell)(v)=a\,\ell(U_a^{-1}v).
\tag{37}
\]

The factor a is carried from the actual second moment representation. It is not set to one or inferred from the cardinality of a base point.

\subsubsection{6. An actual morphism from arithmetic cohomology to its τ-base dual}\label{endpoint-tau_base--an-actual-morphism-from-arithmetic-cohomology-to-its-ux3c4-base-dual}

Put \(g=2\xi\). Retain a finite reflection-stable set Z of actual zeros, with actual multiplicities. Define

\[
A_Z=\bigoplus_{\rho\in Z}\mathcal O_\rho/(g),
\qquad
(J_ZF)_\rho=\sum_{j<m_\rho}
\frac{(\mathcal MF)^{(j)}(\rho)}{j!}z_\rho^j.
\tag{38}
\]

The source identity \(\mathcal M\Theta\phi=gH_\phi\), with \(H_\phi\) entire, proves that the maps

\[
\mathcal T\longrightarrow S_\eta A_Z,
\quad 0\text{ on }+, -,\sigma,\quad J_Z\text{ on }\eta,
\tag{39}
\]

are a sheaf morphism. Applying \(\mathsf A_\tau\) gives precisely the already-specified map \(Q\to A_Z\) in degree one. This finite jet map is onto; the source proof is independence of the finite family \(e^{\rho y}y^j/j!\) on any nonempty interval, followed by a compactly supported smooth right inverse.

On \(A_Z\), the scaling action is

\[
U_a|_\rho=a^\rho\sum_{j<m_\rho}\frac{(\log a)^j}{j!}N_\rho^j,
\quad N_\rho z^j=z^{j+1}.
\tag{40}
\]

Using the reflected germ map

\[
f^\dagger(s)=\overline{f(1-\bar s)},
\quad
\sum c_jz_{\iota\rho}^j\mapsto\sum(-1)^j\bar c_jz_\rho^j,
\tag{41}
\]

retain the perfect finite residue pairing

\[
R_Z(f,h)=\sum_{\rho\in Z}\operatorname{Res}_\rho
\frac{f^\dagger(s)h(s)}{g(s)}\,ds.
\tag{42}
\]

It has \(R_Z(U_af,U_ah)=aR_Z(f,h)\). Its raw skew-Hermitian orientation is retained; no multiplication by a phase is inserted to call it positive.

Precomposing the dual map of (39) with (42) gives an explicit injective linear morphism from the conjugate finite packet:

\[
\boxed{
\overline{A_Z}\longrightarrow
H^{-1}R\operatorname{Hom}_{\mathcal A}(\mathcal T,K_\tau(\mathcal L_1)),
\quad
\bar f\longmapsto
\left([F]\longmapsto\sum_{\rho\in Z}\operatorname{Res}_\rho
\frac{f^\dagger(s)\mathcal MF(s)}{g(s)}\,ds\right).
}
\tag{43}
\]

The target is canonically \(\operatorname{Hom}_E(Q,\mathcal L_1)\) by (28). Well-definedness follows because replacing F by \(F+\Theta\phi\) adds the holomorphic residue expression \(f^\dagger H_\phi\). Injectivity follows from the surjectivity of \(J_Z\) and perfectness of (42). Formula (37) and covariance of (42) prove equivariance.

Thus the duality here is a morphism \textbf{to the right-adjoint object computed on the τ-base}, not merely a pairing written separately from a cohomological model.

The trace contraction is the specified endomorphism

\[
J_g:A_Z\to A_Z,\quad[h]\mapsto[g'h],
\tag{44}
\]

\[
R_Z(f,J_gh)=\sum_\rho m_\rho\overline{f(1-\bar\rho)}h(\rho)
=\operatorname{Tr}(M_{f^\dagger h}|A_Z).
\tag{45}
\]

At a zero of order m, write the actual \(g=u(z)z^m\). Then \(g'\equiv m u(z)z^{m-1}\pmod{z^m}\). The radical of (45) consists of the jets with zero residue value. They map to e under the original internal quotient, while (42) retains their full pairing before that observation.

All maps in (38)--(45) lift to supported fibres. At label A the map is \((A,[F])\mapsto(A,J_ZF)\). A zero functional obtained in (43) or (45) is the supported zero functional at that label, not τ. For inclusions of admitted source labels the resulting squares commute, and the dual arrows are the precomposition arrows (31).

\subsubsection{7. The tensor-amplification operation on the changed base}\label{endpoint-tau_base--the-tensor-amplification-operation-on-the-changed-base}

For the joint arithmetic fibre, take products of the actual chart space P over the underlying point of \(\mathfrak b_\tau\), and external products of its coefficient sheaf. This is the ordinary product of the declared finite topological models with external coefficient tensors. It is not a claim that its structure-sheaf product has already been identified with the schematic square of the full arithmetic space.

The projective resolution for \(P^r\) is the external tensor product of (18). Its differential on homogeneous factors is

\[
d(x_1\otimes\cdots\otimes x_r)=
\sum_{j=1}^r(-1)^{\sum_{k<j}|x_k|}
 x_1\otimes\cdots\otimes dx_j\otimes\cdots\otimes x_r.
\tag{46}
\]

Because the tensor product is over the field E, this yields the explicit Künneth map and, at the top degree,

\[
\boxed{
H^r\mathsf A_{\tau,r}(\mathcal T^{\boxtimes r})\cong Q^{\otimes_E r}.
}
\tag{47}
\]

Indeed the numerator of top cohomology is \(\mathscr B^{\otimes r}\); its relations are the sum of all subspaces obtained by putting \(\Theta V\) in at least one factor. The quotient is exactly the displayed tensor quotient. The joint-support evaluation is a stated functor from the full retained support diagram. Other original support fibres are not removed; their cochain complexes and transports remain in that diagram.

The top-degree kernel relation is again quotiented internally in its supported scalar fibre. Pure external absence is absorbing. A supported tensor whose amplitude becomes a tensor relation maps to the supported zero of that tensor fibre. We do not identify a coefficient external tensor with the Cartesian product of support labels; such a tensor has its specified absorbing basepoint and balancing relations.

The dual right-adjoint object on the product is computed from the dual tensor resolution:

\[
K_{\tau,r}(W)\simeq S_{(\eta,\ldots,\eta)}W[r].
\tag{48}
\]

For \(r=2\), the degree −2 term and both degree −1 terms occur in that resolution with the signs from (46). The shift {[}2{]} is obtained by this actual product calculation, not inserted into (30).

For a residue-value eigenvector at an actual zero \(\rho\), the action in (47) on its r-fold tensor is

\[
a^{r\rho}.
\tag{49}
\]

In particular the exact modulus defect relative to r copies of \(\mathcal L_1^{1/2}\), expressed without rescaling the action, is

\[
\log\frac{|a^{r\rho}|^2}{a^r}
=r(2\Re\rho-1)\log a\quad(a>1).
\tag{50}
\]

This performs the amplification operation on the actual arithmetic spectral quotient over the new base. It does not prove a bound on the left-hand side. Deligne's geometric upper estimate is additional information in his setting; it has not been derived for (47).

\subsubsection{8. What a map to the base does, and what deleting a spectral packet does}\label{endpoint-tau_base--what-a-map-to-the-base-does-and-what-deleting-a-spectral-packet-does}

The structural map (3) sends every point to the base point. The derived pushforward (19) carries its coefficient relations and arithmetic operators along with it. Restriction (20), by contrast, has value the support skeleton on our sheaf. These are two explicitly compared operations on the same model.

For clarity about the proposed ``exceptions land in τ,'' let W and W' be ordinary vector spaces and use their single-active-fibre split lifts. If a \(G(E)\)-linear map

\[
f:W^\tau\to(W')^\tau
\]

sends one supported vector v to τ, then it sends every supported vector to τ. Proof: \(f(0_W^\bullet)=f(ev)=ef(v)=\tau\). For every w, \(ef(w)=f(ew)=\tau\). In \((W')^\tau\), the only element whose e-multiple is τ is τ itself. This result is for this specified target; it is not asserted for arbitrary semimodules with a nontrivial bottom group fibre.

A labelled projection can delete a whole packet while retaining other labels, but it has a calculable cohomological contribution. For a reflection-stable finite set Z split into C and B, the actual projection is

\[
A_Z\twoheadrightarrow A_C,\qquad\ker=A_B.
\tag{51}
\]

The difference of convolution traces is

\[
\operatorname{Tr}(H_h|A_Z)-\operatorname{Tr}(H_h|A_C)
=\sum_{\rho\in B}m_\rho\mathcal Mh(\rho).
\tag{52}
\]

This is nonzero for some admissible compactly supported smooth h whenever B is nonempty, by finite Mellin interpolation. Its split value for that test is a supported nonzero number. The cone of the projection \(A_Z[-1]\to A_C[-1]\) has \(H^0=A_B\); that is the precise cohomological trace lost by deletion.

Consequently a structural projection of the entire geometry to the τ-base does not delete a packet. A packet deletion is (51), with the separate kernel and trace (52). Establishing that all actual exceptional contributions are absent requires this kernel or its full trace contribution to vanish by arithmetic geometry, not merely by relabelling its support.

\subsubsection{9. Result and next mathematically defined target}\label{endpoint-tau_base--result-and-next-mathematically-defined-target}

The completed construction in this note is

\[
\widetilde{\mathcal T}\text{ over }\mathfrak b_\tau
\ \xrightarrow{\ \mathsf A_\tau\ }\
\widetilde H^1_\tau
\ \xrightarrow{\ J_Z\ }\
A_Z^\tau,
\]

with the computed adjoint object \(K_\tau(\mathcal L_1)\), the dual morphism (43), the trace contraction (45), and all tensor powers (47). It is a characteristic-zero scaling calculation, not a purity calculation on a finite-field curve.

The arithmetic quotient \(G(\mathbb Z)\to\mathbb Z\) is kept in the structural map (8). The current theta chart model sees its generic arithmetic image; a global prime-sensitive geometric realization with a Lefschetz theorem and two-sided weight estimates is not obtained from this finite chart topology alone. The specific remaining quantity in this model is the tensor-growth defect (50), or equivalently the sign of the full trace pairing (45) after its global arithmetic realization. Neither is set to zero in the note.

The new work is the actual change of base, the computed right adjoint, and the cohomological placement of the finite arithmetic duality maps on that base. It narrows the task to a property of these explicit objects rather than repeating the assertion that a shared support should imply purity.

\subsubsection{Sources read for this calculation}\label{endpoint-tau_base--sources-read-for-this-calculation}

\begin{enumerate}
\def\labelenumi{\arabic{enumi}.}
\tightlist
\item
  \nolinkurl{globalization_note.tex}: lines 432--563 (localizations and stalks), 673--751 (contracted monoid algebra and blueprint remark). The separate supplied Lean-summary note provides the correction \texttt{{[}tau{]}=0} and the implemented support-fibre operations. Those source results are not claimed as new.
\item
  Retained \nolinkurl{Split_Support_Actual_Comparison/CONTINUATION.tex}: test spaces, Fourier and theta maps, the arithmetic quotient and finite jet realization. This is an earlier research derivation, not an independent verification of every global statement in its package.
\item
  The supplied French Weil II TeX excerpts, original §6.2.1--6.2.7 and §3.3: the right-adjoint definition of a dualizing complex, opposite dual-weight bounds, proper pushforward, and the arithmetic-base twist/degree equation. The excerpts' historical-transcription status remains recorded. Reading these windows is not a claim to read or audit the whole paper in this turn.
\item
  O. Lorscheid, \emph{The geometry of blueprints. Part I} (arXiv:1103.1745): external framework for the declared blueprint base. The finite-chart model and adjoint calculation above are proved here; they are not attributed to that paper.
\end{enumerate}

\subsubsection{Checks}\label{endpoint-tau_base--checks}

The accompanying executable checks finite instances of the stated poset resolution, dual complex, internal support differential, action compatibility, tensor differentials and trace deletion. They use explicit exceptions rather than optimizable \nolinkurl{assert} statements. They are regression checks of these identities, not a formal proof of RH, of Deligne's theorem, or of a global arithmetic realization.
